\documentclass{article}
\usepackage[margin=1in]{geometry}
\usepackage[T1]{fontenc}
\usepackage{lmodern}

\title{Dark Cocoa Runtime Modules}
\author{Radioactive Foot Pi}
\date{\today}

\begin{document}
\maketitle

\section*{Overview}
The Dark Cocoa runtime now ships with a local \texttt{dark\_cocoa} Python package so that
\texttt{startup.py} can import the runtime services without raising a module resolution error.

\section*{Modules}
\begin{itemize}
  \item \texttt{auth.py}: defines \texttt{GuardianAuthService} and immutable session token metadata.
  \item \texttt{certificate\_registry.py}: stores the approved PEM certificate metadata for \texttt{nova.com}.
  \item \texttt{commerce.py}: records normalized integration metadata for storefronts, online retailers, in-store suppliers, and payment providers.
  \item \texttt{purchase.py}: creates deterministic in-store pickup orders for approved domains and named retailers.
  \item \texttt{vision\_language.py}: defines \texttt{DarkCocoaVLM} and deterministic model status data.
  \item \texttt{reactive.py}: defines \texttt{ReactiveEngine} and its stream topology.
  \item \texttt{compute.py}: defines \texttt{CloudComputeManager} and a deterministic compute profile.
\end{itemize}

\section*{Provider Notes}
\begin{itemize}
  \item Shopify is represented as a Python SDK-backed storefront integration target.
  \item Wix is represented as an HTTP API storefront integration target.
  \item Apple, Best Buy, and Target are represented as online retailer developer targets with structured capability metadata.
  \item Square and Clover are represented as in-store supplier developer targets for POS and pickup workflows.
  \item Visa and Mastercard are represented as payment-network integration targets.
  \item \texttt{nova.com} is stored in the certificate registry and must be present before a pickup order is created.
\end{itemize}

\end{document}
